\chapter{Foot Drop}

\section*{Definition}
Foot drop is the inability to dorsiflex the ankle, resulting in a steppage gait where you lifts the knee higher to clear the foot during the swing phase of gait. It is caused by weakness of the tibialis anterior muscle and may be due to upper motor neuron, lower motor neuron, or peripheral nerve pathology.

\section*{What Happens in Your Body}
Dorsiflexion of the ankle is primarily mediated by the tibialis anterior muscle, innervated by the deep fibular (peroneal) nerve (L4--S1 nerve roots via the sciatic nerve). Lesions at any point along this pathway cause foot drop. Lower motor neuron (LMN) lesions produce flaccid weakness, atrophy, and fasciculations. Upper motor neuron (UMN) lesions (stroke, multiple sclerosis, spinal cord injury) produce spastic weakness with hyperreflexia and an extensor plantar response. The common fibular nerve is vulnerable to compression at the fibular neck due to its superficial location and poor vascular supply.

\section*{Causes}
\begin{itemize}
  \item \textbf{Peripheral nerve:} Common fibular (peroneal) neuropathy at the fibular neck (compression, knee crossing, casting, prolonged squatting, weight loss), sciatic neuropathy, L5 radiculopathy
  \item \textbf{Spinal:} Lumbar disc herniation (L4--L5), spinal stenosis, cauda equina syndrome, trauma
  \item \textbf{Motor neuron:} Amyotrophic lateral sclerosis, spinal muscular atrophy, polio
  \item \textbf{Neuromuscular junction:} Myasthenia gravis (variable, fatigable)
  \item \textbf{Muscle:} Tibialis anterior rupture, myopathy (inclusion body myositis, distal myopathies)
  \item \textbf{Upper motor neuron:} Stroke, multiple sclerosis, cerebral palsy, brain tumor
  \item \textbf{Systemic:} Charcot-Marie-Tooth disease, hereditary neuropathy with liability to pressure palsy (HNPP), diabetes, vasculitic neuropathy
  \item \textbf{Iatrogenic:} Post-hip arthroplasty, post-knee arthroplasty, post-surgical positioning
\end{itemize}

\section*{When to See a Doctor}
See your doctor if:
\begin{itemize}
  \item Acute onset of foot drop with back pain, saddle anesthesia, or urinary retention (rule out cauda equina syndrome)
  \item Foot drop after trauma or surgery
  \item Progressive foot drop with weakness in other muscle groups
  \item Foot drop with gait impairment causing falls
\end{itemize}

\section*{Related Symptoms}
\begin{itemize}
  \item Slapping gait (foot slaps the ground during stance phase)
  \item High-stepping gait (excessive hip and knee flexion to clear the foot)
  \item Numbness or paresthesias over the dorsum of the foot and lateral leg
  \item Weakness of ankle eversion (fibularis muscles) in common fibular neuropathy
  \item Inversion of the foot at rest (unopposed tibialis posterior)
\end{itemize}
\textbf{Important:} Acute foot drop with severe low back pain and urinary retention is cauda equina syndrome, a surgical emergency requiring immediate spinal imaging.

\section*{Self-Care / Home Management}
\begin{itemize}
  \item Ankle-foot orthosis (AFO) to maintain dorsiflexion during gait
  \item Remove tripping hazards in the home environment
  \item Avoid crossing legs at the knee to prevent peroneal compression
  \item Over-the-counter NSAIDs for pain if associated with radiculopathy
  \item Skin inspection of the foot and ankle due to altered sensation
\end{itemize}

\section*{How Your Doctor Will Evaluate You}
\begin{itemize}
  \item Nerve conduction studies and electromyography to localize the lesion (fibular nerve, sciatic, L5 root, motor neuron)
  \item MRI of the lumbar spine if L4--L5 radiculopathy is suspected
  \item MRI or ultrasound of the knee/common fibular nerve for compression or mass lesion
  \item MRI brain or cervical spine if UMN signs suggest stroke, MS, or cervical myelopathy
  \item Laboratory evaluation: HbA1c, vitamin B12, TSH, ESR, ANA, anti-MAG, genetic testing for HNPP/CMT
\end{itemize}

\section*{Treatment Options}
\begin{itemize}
  \item Ankle-foot orthosis (AFO) for immediate gait stability and fall prevention
  \item Physical therapy with ankle dorsiflexion strengthening and gait retraining
  \item Surgical decompression of the common fibular nerve at the fibular neck for compressive neuropathy
  \item Microsurgical repair or nerve grafting for nerve transection
  \item Posterior tibial tendon transfer for permanent foot drop
  \item Functional electrical stimulation (FES) of the common fibular nerve during the swing phase
  \item Treatment of underlying cause (discectomy for herniation, disease-modifying therapy for MS, glycemic control for diabetes)
\end{itemize}

\section*{References}
\begin{thebibliography}{99}
\bibitem{ref1} Stewart JD. "Foot Drop: Where, Why and What to Do?" Pract Neurol. 2008;8(3):158--169.
\bibitem{ref2} Aprile I, Padua L, Padua R, et al. "Peroneal Nerve Palsy: A Clinical and Electrophysiological Study." J Foot Ankle Surg. 2000;39(3):176--179.
\bibitem{ref3} Van Langenhove M, Pollefliet A, Vandenberghe R. "Drop Foot." BMJ. 2013;346:f2187.
\bibitem{ref4} Masakado Y, Kawakami M, Suzuki K, et al. "The Pathophysiology and Management of Foot Drop." J Rehabil Med. 2018;50(7):583--590.
\bibitem{ref5} Preston DC, Shapiro BE. "Electromyography and Neuromuscular Disorders." 4th ed. Elsevier; 2021.
\end{thebibliography}
